\section{Motivating Examples}

\begin{figure}[!t]
    \begin{minted}[xleftmargin=5pt, numbersep=4pt, linenos = true, fontsize = \footnotesize, escapeinside=??]{OCaml}
    let rec merge (l1: int list) (l2: int list) : int list =
      match l1, l2 with
      | [], l2 -> l2
      | l1, [] -> l1
      | h1 :: t1, h2 :: t2 ->
        if h1 < h2 then h1 :: h2 :: merge t1 t2
        else h2 :: h1 :: merge t1 t2
    
    let rec merge (l1: int list) (l2: int list) : int list =
        match l1, l2 with
        | [], l2 -> l2
        | l1, [] -> l1
        | h1 :: t1, h2 :: t2 ->
          if h1 < h2 then h1 :: merge t1 (h2 :: t2)
          else h2 :: merge (h1 :: t1) t2

    let sorted (l: int list) : bool =
      match l with
      | [] | [_] -> true
      | h1 :: t ->
        h1 <= (hd t) && sorted t
    \end{minted}
    \caption{
    The $\Code{merge}$ function is supposed to merge two sorted lists. If either of the input lists is empty, \Code{merge} returns the other (lines $3$--$4$). Otherwise, the function merges the lists recursively (lines $6$--$8$). However, the first implementation is buggy: in the case $\Code{h1} < \Code{h2}$, it places $\Code{h2}$ before the result of merging $\Code{t1}$ and $\Code{t2}$, incorrectly assuming that $\Code{h2}$ is smaller than all elements in $\Code{t1}$. }
    \label{fig:motivation}
\end{figure}

\paragraph{Verification Failure} We attempt to prove the \Code{merge} function preserves the sorting property, that is,
\begin{align*}
    \forall l_1\;l_2.\ \I{sorted}(l_1) \impl \I{sorted}(l_2) \impl \I{sorted}(\Code{merge}(l_1, l_2))
\end{align*}\noindent
or, equivalently, the following refinement type:
\begin{align*}
    \Code{merge} : \ort{il}{\I{sorted}(\vnu)} \sarr \ort{il}{\I{sorted}(\vnu)} \sarr \ort{il}{\I{sorted}(\vnu)}
\end{align*}\noindent
However, neither of the two implementations is typeable. The type check fails at the $\Code{::}$ constructor. For example, on line $14$ (the correct implementation), we cannot prove that the element $\Code{h1}$ is smaller than all elements in the result of the recursive call, whose type is just $\ort{il}{\I{sorted}(\vnu)}$. The first (buggy) implementation has the same problem, and we cannot distinguish between them.

\paragraph{More Precise Type} This problem can be solved by providing a more precise type for the $\Code{merge}$ function. For example, we can additionally consider the membership.
\begin{align*}
    \Code{merge} : &l_1{:}\ort{il}{\I{sorted}(\vnu)} \sarr l_2{:}\ort{il}{\I{sorted}(\vnu)} \sarr 
    \\&\ort{il}{\I{sorted}(\vnu) \land \I{elemSet}(l_1) \cup \I{elemSet}(l_2) = \I{elemSet}(\vnu)}
\end{align*}\noindent
where $\I{elemSet}(l)$ is the set of elements in the list $l$. However, coming up with this new predicate (measure) is not easy, and the relation between $\I{sorted}$ and $\I{elemSet}$ still requires additional lemmas (e.g., $\I{sorted}(x::l) \impl \forall y. y \in \I{elemSet}(l) \impl x \leq y$).

\paragraph{Under-Approximation Type Shows Incorrectness} The first implementation is indeed buggy. One possible counterexample is $\Code{merge\;[3]\;[1;2] = [1;3;2]}$. We can use under-approximation reasoning to find a more general erroneous condition.
\begin{align*}
    \Code{merge}& : \ort{il}{\I{sorted}(\vnu) \land \vnu \not= [] } \sarr \urt{il}{\I{sorted}(\vnu)} \sarr \urt{il}{\neg \I{sorted}(\vnu)}
    \\\Code{merge}& : \ort{il}{\vnu = [] } \sarr y{:}\ort{il}{\I{sorted}(\vnu)} \sarr \urt{il}{\vnu = y}
    \\\Code{merge}& : x{:}\ort{il}{\I{sorted}(\vnu)} \sarr \ort{il}{\vnu = []} \sarr \urt{il}{\vnu = x}
\end{align*}\noindent
If the first list is not empty, there must exist a sorted list that makes the merge result unsorted. To prove the type, we need two overapproximated parameter types (shown above).

\paragraph{Under-Approximation Helps Refine Correctness Specification} The second implementation is indeed correct, and under-approximation will fail to prove the reachability of any potential buggy path. For example, it fails to prove that $\Code{h1\;::\;merge\;t1\;(h2\;::\;t2)}$ is unsorted, which means that the list $l$ that makes $\Code{h1\;::l}$ unsorted is not reachable.
\begin{align*}
    \Code{merge}& : \Code{h1} \garr \ort{il}{\I{sorted}(\vnu) \land \Code{h1} \leq \I{hd}(\vnu)} \sarr \ort{il}{\I{sorted}(\vnu) \land \Code{h1} \leq \I{hd}(\vnu)} 
    \\&\sarr \Code{Unreachable}(\urt{il}{\neg \I{sorted}(\Code{h1}::\vnu)})
\end{align*}\noindent
Note that $\Code{h1}$ is less than the head of the first list (according to the definition of $\Code{sorted}$ on line $21$), and also less than the head of the second list (since the head of the second list is less than the head of the first list). Then, we can use this fact to refine the original specification.
\begin{align*}
    \Code{merge}& : \Code{h1} \garr \ort{il}{\I{sorted}(\vnu) \land \Code{h1} \leq \I{hd}(\vnu)} \sarr \ort{il}{\I{sorted}(\vnu) \land \Code{h1} \leq \I{hd}(\vnu)} 
    \\&\sarr \ort{il}{\I{sorted}(\Code{h1}::\vnu)}
\end{align*}\noindent
This refined specification is sufficient to make the type check succeed. For example, consider the branch on line $14$. We have:
\begin{align*}
    &\Code{h1}{:}\Int, \Code{h2}{:}\ort{\Int}{\Code{h1} \leq \vnu}, \Code{t1}{:}\ort{il}{\I{sorted}(\vnu) \land \Code{h1} \leq \I{hd}(\vnu)}, ...\vdash 
    \\&\quad \Code{merge}\;\Code{t1}\;(\Code{h2}::\Code{t2}) : \ort{il}{\I{sorted}(\Code{h1}::\vnu)}
\end{align*}\noindent
Note that $\Code{h1} < \I{hd}(\vnu)$ and $\Code{h1} < \Code{h2} = \I{hd}(\Code{h2}::\Code{t2})$.
On the other hand, for an arbitrary ghost variable $i$ that less than head elements of the two lists,
\begin{align*}
    &i{:}\ort{\Int}{\vnu \leq \Code{h1} \land \vnu \leq \Code{h2}}, \Code{h1}{:}\Int, ...\vdash 
    \\&\quad (\Code{h1} :: \Code{merge}\;\Code{t1}\;(\Code{h2}::\Code{t2})) : \ort{il}{\I{sorted}(i::\Code{h1}::...)}
\end{align*}\noindent
This can be proved from the definition of $\Code{sorted}$ by the facts that $\Code{h1}::...$ is sorted (proved above), and $i$ is less than $\I{hd}(\Code{h1}::...) = \Code{h1}$.

\paragraph{Discussion} Note that the refined specification is much weaker than the one considering membership; also, we do not add any new predicates (e.g., $\I{elemSet}$) to the specification, and no new lemmas are required.

\section{Old: Motivating Examples}

\begin{figure}[!t]
    \begin{minted}[xleftmargin=5pt, numbersep=4pt, linenos = true, fontsize = \footnotesize, escapeinside=??]{OCaml}
    let rec merge (l1: int list) (l2: int list) : int list =
      match l1, l2 with
      | [], l2 -> l2
      | l1, [] -> l1
      | h1 :: t1, h2 :: t2 ->
        if h1 < h2 then h1 :: h2 :: merge t1 t2
        else if h1 > h2 then h2 :: h1 :: merge t1 t2
        else h1 :: h2 :: merge t1 t2
    
    let rec mergesort (split: int list -> int list * int list) (l: int list) =
        match l with
        | [] | [_] -> l
        | _ ->
            let l1, l2 = split l in
            merge (mergesort l1) (mergesort l2)
    \end{minted}
    \caption{\Code{mergesort} implements the merge sort algorithm: it splits the list, recursively sorts the sublists, and merges the results (lines $14$--$15$). \Code{mergesort} is parameterized by the input \Code{split} function, which can use different splitting strategies.
    The $\Code{merge}$ function is supposed to merge two sorted lists. If either of the input lists is empty, \Code{merge} returns the other (lines $3$--$4$). Otherwise, the function merges the lists recursively (lines $6$--$8$). However, this implementation is buggy: in the case $\Code{h1} < \Code{h2}$, it places $\Code{h2}$ before the result of merging $\Code{t1}$ and $\Code{t2}$, incorrectly assuming that $\Code{h2}$ is smaller than all elements in $\Code{t1}$. }
    \label{fig:motivation}
\end{figure}

\paragraph{Proving Incorrectness} We attempt to prove the incorrectness of \Code{mergesort}, as shown in Figure \ref{fig:motivation}; that is, it violates the sorting property.
\begin{align*}
    \text{sorting property: } & \forall l\;l'.\ l' = \Code{mergesort}\ l \impl \I{sorted}(l')
    \\ \text{incorrectness: } & \exists l\;l'.\ l' = \Code{mergesort}\ l \land \neg \I{sorted}(l')
\end{align*}\noindent
More precisely, we would like to know under what conditions (e.g., shape of input list, properties of the \Code{split} function) the \Code{mergesort} function is incorrect.

\paragraph{Example} For input $\Code{l1 = [1;2],\ l2 = [5;6]}$, the \Code{merge} function will produce the unsorted output $\Code{[1;5;2;6]}$. If we choose the second implementation of the \Code{split} function shown in Figure \ref{fig:split}, the \Code{mergesort} function will produce the unsorted list from the input $\Code{[2;1;6;5]}$.
\begin{figure}[!t]
    \begin{minted}[xleftmargin=5pt, numbersep=4pt, linenos = true, fontsize = \footnotesize, escapeinside=??]{OCaml}
    (* the first list will not be less than the second list *)
    let split (l: int list) : int list * int list =
        let rec aux l1 l2 = function
          | [] -> (l1, l2)
          | [x] -> (x :: l1, l2)
          | x :: y :: rest ->
              aux (x :: l1) (y :: l2) rest
        in
        aux [] [] l

    (* the first list will be no longer than the second list *)
    let split (l: int list) : int list * int list =
        let n = List.length l in
        let n1 = n / 2 in
        let n2 = n - n1 in
        (List.sublist l 0 n1, List.sublist l n1 n)

    (* random pick an element as pivot and split the list into two parts ... *)
    \end{minted}
    \caption{Two possible implementations of the \Code{split} function.}
    \label{fig:split}
\end{figure}

\paragraph{Challenges} There are several challenges in proving incorrectness.
\begin{enumerate}
    \item Higher-order functions, which prevent traditional testing and symbolic execution approaches. We also want to reject some trivial invalid input functions.
    \item Compositional reasoning: we would like to find a precise incorrectness precondition for \Code{mergesort} to enable further proofs over different implementations of \Code{split}. For example, the following specification is not allowed:
    \begin{align*}
        \Code{split} : \ort{\List[\Int]}{\top} \sarr \urt{\List[\Int]}{\top} \times \urt{\List[\Int]}{\top}
    \end{align*}\noindent
    where $\Code{split}$ should follow some basic properties, e.g., preserve membership of the input list.
    Then the incorrectness precondition of $\Code{merge}$ should also be precise; i.e., it cannot just be 
    \begin{align*}
        \Code{merge} : \urt{\List[\Int]}{\top} \sarr \urt{\List[\Int]}{\top} \sarr \Code{unsorted}
    \end{align*}\noindent
    \item Recursive functions, where the incorrectness execution path involves multiple recursive calls. For some calls, we need a correctness specification; for others, we need an incorrectness specification.
\end{enumerate}

\paragraph{Solution: Using Over- and Under- Approximation}
\begin{enumerate}
    \item Higher-order function: We provide an ``over-style'' specification of the input function:
    \begin{align*}
        &\Code{mergesort}:
        \\&\quad \ol{\Code{split}}{:}(\ol{l}{:}\rt{il}{\top} \sarr \ol{\Code{l1}}{:}\rt{il}{\top} \times \rt{il}{\I{len}(\vnu) = \I{len}(\Code{l1}) + \I{len}(\vnu) \land \I{mem}\ldots}) \sarr
        \\&\quad \ol{\Code{l}}{:}\rt{il}{\top} \sarr
        \\&\quad \rt{il}{\I{sorted}(\vnu)}
    \end{align*}\noindent
    Then we ``negate'' it to get the ``under-style'' specification:
    \begin{align*}
        &\Code{mergesort}:
        \\&\quad \ul{\Code{split}}{:}(\ol{l}{:}\rt{il}{\top} \sarr \ol{\Code{l1}}{:}\rt{il}{\top} \times \rt{il}{\I{len}(\vnu) = \I{len}(\Code{l1}) + \I{len}(\vnu) \land \I{mem}\ldots}) \sarr
        \\&\quad \ul{\Code{l}}{:}\rt{il}{\top} \sarr
        \\&\quad \rt{il}{\neg \I{sorted}(\vnu)}
    \end{align*}\noindent
    Note that the function type of $\Code{split}$ is not flipped (it is still an over-style specification), but the way we treat the parameter $\Code{split}$ is changed to under-style (similar to higher-order model checking). Our proof should show that the program can be typechecked against some supertype of this type. Similarly, the $\Code{merge}$ function should have an over-style specification:
    \begin{align*}
        &\Code{merge}: \ol{\Code{l1}}{:}\rt{il}{\I{sorted}(\vnu)} \sarr \ol{\Code{l2}}{:}\rt{il}{\I{sorted}(\vnu)} \sarr \rt{il}{\top}
    \end{align*}\noindent
    \item Mixed specification: The $\Code{merge}$ function requires two parameters; if both of them are underapproximated, it is too specific and hard to match the actual implementation. For example, the two implementations of the $\Code{split}$ function shown in Figure \ref{fig:split} have different guarantees on the lengths of the output lists. However, the following specification is more general:
    \begin{align*}
        &\Code{merge}:
        \\&\quad \ol{\Code{l1}}{:}\rt{il}{\I{len}(\vnu) \geq 2 \land \I{sorted}(\vnu)} \sarr \ul{\Code{l2}}{:}\rt{il}{\I{len}(\vnu) = 2 \land \I{sorted}(\vnu)} \sarr \rt{il}{\neg \I{sorted}(\vnu)}
        \\&\quad \ol{\Code{l1}}{:}\rt{il}{\I{len}(\vnu) \geq 2 \land \I{sorted}(\vnu)} \sarr \ul{\Code{l2}}{:}\unionty_{i{:}\Int}\rt{il}{\I{len}(\vnu) = i \land i \geq 2 \land \I{sorted}(\vnu)} \sarr \rt{il}{\neg \I{sorted}(\vnu)}
    \end{align*}\noindent
    This can lead to a more general specification of the $\Code{split}$ function (i.e., the incorrectness precondition):
    \begin{align*}
        &\Code{split}:
        \\&\quad \ul{l}{:}\rt{il}{\I{len}(\vnu) \geq 4} \sarr
        \\&\quad \ol{\Code{l1}}{:}\rt{il}{\I{len}(\vnu) \geq 2} \times \urt{il}{\I{len}(\vnu) = 2}
    \end{align*}\noindent
    This type can fit both implementations of the $\Code{split}$ function shown in Figure \ref{fig:split}.
    \item Intersection type: The recursive function $\Code{merge}$ also requires correctness specifications (otherwise, recursive calls cannot be typechecked):
    \begin{align*}
        &\Code{merge}:
        \\&\quad \ol{\Code{l1}}{:}\rt{il}{\I{len}(\vnu) < 2 \land \I{sorted}(\vnu)} \sarr \ol{\Code{l2}}{:}\rt{il}{\I{len}(\vnu) < 2 \land \I{sorted}(\vnu)} \sarr 
        \\&\quad \rt{il}{\I{sorted}(\vnu) \land \I{len}(\vnu) = \I{len}(\Code{l1}) + \I{len}(\Code{l2}) \land \I{mem}\ldots}
    \end{align*}\noindent
    Thus, the $\Code{merge}$ function should be annotated with intersected types.
\end{enumerate}